\chapter{Vectores y rectas en el plano}\label{ch:g11:vect}

Los \omterm{def:g11:vect:vector}{vectores} codifican desplazamientos; las rectas son sus trayectorias. La
herramienta nueva y clave de este capítulo es el \emph{\omterm{def:g11:vect:det}{determinante}}, un
único número que decide si dos \omterm{def:g11:vect:vector}{vectores} son \omterm{def:g11:vect:collinear}{colineales} y que, a partir de
ahí, produce ecuaciones de rectas y resuelve por puro cálculo las
cuestiones de alineación y paralelismo. Las mismas ideas se extienden al
espacio en el \cref{ch:g12:space}.

\section{Vectores: repaso}

\begin{definition}[Vector, coordenadas]\label{def:g11:vect:vector}
Un \emph{vector}\index{vector} $\vec u$ representa un desplazamiento:
todas las flechas de la misma dirección, sentido y longitud representan el
mismo vector. En un \omterm{def:g10:coordgeom:system}{sistema de coordenadas}, $\vec u\,(x, y)$ recoge el
desplazamiento ($x$ en horizontal, $y$ en vertical); para dos puntos
$A(x_A, y_A)$ y $B(x_B, y_B)$,
\[
\vect{AB}\,\bigl(x_B - x_A,\ y_B - y_A\bigr).
\]
Los vectores se suman coordenada a coordenada, y $\lambda\vec u$ tiene
\omterm{def:g10:coordgeom:system}{coordenadas} $(\lambda x, \lambda y)$.
\end{definition}

\begin{proposition}[Punto medio]\label{prop:g11:vect:midpoint}
El \omterm{prop:g10:coordgeom:midpoint}{punto medio} $M$ del segmento $[AB]$ tiene \omterm{def:g10:coordgeom:system}{coordenadas}
$\left(\frac{x_A + x_B}{2},\ \frac{y_A + y_B}{2}\right)$.
\end{proposition}

\begin{proof}
$M$ está caracterizado por $\vect{AM} = \frac12\vect{AB}$: coordenada a
coordenada, $x_M - x_A = \frac12(x_B - x_A)$, luego
$x_M = \frac{x_A + x_B}{2}$, y análogamente para $y_M$.
\end{proof}

\section{Colinealidad y determinante}

\begin{definition}[Vectores colineales]\label{def:g11:vect:collinear}
Dos \omterm{def:g11:vect:vector}{vectores} son \emph{colineales}\index{vectores colineales} si uno es
múltiplo del otro: $\vec v = \lambda \vec u$ para algún $\lambda \in \R$
(o bien $\vec u = \vec 0$). Geométricamente: tienen la misma dirección, o
uno de ellos es nulo.
\end{definition}

\begin{definition}[Determinante]\label{def:g11:vect:det}
El \emph{determinante}\index{determinante} de $\vec u\,(x, y)$ y
$\vec v\,(x', y')$ es
\[
\det(\vec u, \vec v) =
\begin{vmatrix} x & x' \\ y & y' \end{vmatrix}
= x y' - x' y .
\]
\end{definition}

\begin{theorem}[Criterio de colinealidad]\label{thm:g11:vect:collinearity}
Dos \omterm{def:g11:vect:vector}{vectores} $\vec u$ y $\vec v$ son \omterm{def:g11:vect:collinear}{colineales} si y solo si
$\det(\vec u, \vec v) = 0$.
\end{theorem}

\begin{proof}
($\Rightarrow$) Si $\vec v = \lambda\vec u$, entonces $x' = \lambda x$ e
$y' = \lambda y$, luego
$\det(\vec u, \vec v) = x(\lambda y) - (\lambda x)y = 0$; si
$\vec u = \vec 0$, el \omterm{def:g11:vect:det}{determinante} es trivialmente $0$.

($\Leftarrow$) Supongamos $xy' - x'y = 0$ con $\vec u \neq \vec 0$,
digamos $x \neq 0$ (el caso $y \neq 0$ es simétrico). Tomamos
$\lambda = \frac{x'}{x}$. Entonces $x' = \lambda x$ por construcción, y la
relación $xy' = x'y = \lambda x y$ da, dividiendo entre $x \neq 0$,
$y' = \lambda y$. Por tanto, $\vec v = \lambda \vec u$.
\end{proof}

\begin{example}\label{ex:g11:vect:collinear}
$\vec u\,(3, -2)$ y $\vec v\,(-6, 4)$:
$\det = 3 \times 4 - (-6) \times (-2) = 12 - 12 = 0$, \omterm{def:g11:vect:collinear}{colineales} (en
efecto, $\vec v = -2\vec u$). $\vec u\,(3, -2)$ y $\vec w\,(2, 1)$:
$\det = 3 + 4 = 7 \neq 0$, no \omterm{def:g11:vect:collinear}{colineales}.
\end{example}

\begin{method}[Alineación de tres puntos]\label{met:g11:vect:alignment}
$A$, $B$, $C$ están alineados si y solo si $\vect{AB}$ y $\vect{AC}$ son
\omterm{def:g11:vect:collinear}{colineales}: calcula los dos \omterm{def:g11:vect:vector}{vectores} y comprueba que
$\det\bigl(\vect{AB}, \vect{AC}\bigr) = 0$.
\end{method}

\section{Ecuaciones generales de rectas}

\begin{definition}[Vector director]\label{def:g11:vect:direction}
Un \emph{vector director}\index{vector director} de una recta $d$ es
cualquier \omterm{def:g11:vect:vector}{vector} no nulo $\vec u$ tal que $\vect{AB}$ es colineal con
$\vec u$ para todos los puntos $A, B$ de $d$.
\end{definition}

\begin{theorem}[Ecuación general]\label{thm:g11:vect:cartesian}
Una recta con \omterm{def:g11:vect:direction}{vector director} $\vec u\,(-b, a)$ admite una \omterm{def:g10:algebra:equation}{ecuación} de la
forma
\[
ax + by + c = 0
\qquad (a, b) \neq (0, 0).
\]
Recíprocamente, toda \omterm{def:g10:algebra:equation}{ecuación} de ese tipo describe una recta con \omterm{def:g11:vect:direction}{vector
director} $\vec u\,(-b, a)$.
\end{theorem}

\begin{proof}
Sea $d$ la recta que pasa por $A(x_A, y_A)$ con dirección
$\vec u\,(-b, a)$. Un punto $M(x, y)$ está en $d$ exactamente cuando
$\vect{AM}$ es colineal con $\vec u$, es decir
(\cref{thm:g11:vect:collinearity}), cuando
\[
0 = \det\bigl(\vect{AM}, \vec u\bigr)
= (x - x_A)\,a - (-b)(y - y_A)
= ax + by - (a x_A + b y_A),
\]
una \omterm{def:g10:algebra:equation}{ecuación} de la forma anunciada con $c = -(a x_A + b y_A)$.
Recíprocamente, las soluciones de $ax + by + c = 0$ con, digamos,
$b \neq 0$ son los puntos $\left(x, -\frac{a x + c}{b}\right)$: restar dos
de ellos muestra que toda cuerda es colineal con $(-b, a)$, y siempre
existe alguna solución; se trata de la recta que pasa por ella con
dirección $(-b, a)$.
\end{proof}

\begin{method}[Recta que pasa por dos puntos]\label{met:g11:vect:twopoints}
Para hallar una \omterm{def:g10:algebra:equation}{ecuación} de la recta $(AB)$: calcula el \omterm{def:g11:vect:direction}{vector director}
$\vect{AB}\,(x_B - x_A, y_B - y_A)$; identifícalo con $(-b, a)$, es decir,
toma $a = y_B - y_A$ y $b = -(x_B - x_A)$; y determina $c$ sustituyendo
las \omterm{def:g10:coordgeom:system}{coordenadas} de $A$.
\end{method}

\begin{example}\label{ex:g11:vect:twopoints}
Recta que pasa por $A(1, 2)$ y $B(4, 3)$: $\vect{AB}\,(3, 1)$, luego
$a = 1$, $b = -3$, y la \omterm{def:g10:algebra:equation}{ecuación} es $x - 3y + c = 0$. Sustituyendo $A$:
$1 - 6 + c = 0$, luego $c = 5$. \omterm{def:g10:algebra:equation}{Ecuación}: $x - 3y + 5 = 0$. Comprobación
con $B$: $4 - 9 + 5 = 0$.
\end{example}

\begin{remark}
Cuando $b \neq 0$, la \omterm{def:g10:algebra:equation}{ecuación} se puede despejar en $y$: $y = mx + p$ con
\omterm{def:g10:reffunc:affine}{pendiente} $m = -\frac ab$. La forma general $ax + by + c = 0$ es más
amplia: también cubre las rectas verticales ($b = 0$), que no tienen
\omterm{def:g10:reffunc:affine}{pendiente}.
\end{remark}

\begin{omfigure}
\begin{tikzpicture}
\begin{axis}[omaxis,
  width=8.6cm, height=5.8cm,
  xmin=-1.6, xmax=5.6, ymin=-0.9, ymax=4.4,
  xtick={1,4}, ytick={2,3}]
  \addplot[omDef, thick, domain=-1.4:5.4] {(x+5)/3};
  \fill (axis cs:1,2) circle (1.5pt)
    node[above left, font=\small] {$A$};
  \fill (axis cs:4,3) circle (1.5pt)
    node[above left, font=\small] {$B$};
  \draw[-Stealth, omThm, thick] (axis cs:1,2) -- (axis cs:2.5,2.5)
    node[midway, below right, font=\small] {$\vec u$};
\end{axis}
\end{tikzpicture}

{\small La recta $x - 3y + 5 = 0$ del \cref{ex:g11:vect:twopoints}, su
\omterm{def:g11:vect:direction}{vector director} $\vec u = \vect{AB}$ (en rojo, aquí a escala $\frac12$) y
los dos puntos usados para construir la \omterm{def:g10:algebra:equation}{ecuación}.}
\end{omfigure}

\section{Representación paramétrica}

\begin{definition}[Representación paramétrica]\label{def:g11:vect:parametric}
La recta que pasa por $A(x_A, y_A)$ con dirección
$\vec u\,(\alpha, \beta)$ es el conjunto de los puntos
\[
M\bigl(x_A + t\alpha,\ y_A + t\beta\bigr), \qquad t \in \R .
\]
El número $t$ es el \emph{parámetro}: cada valor de $t$ produce un punto
de la recta, como si se recorriera a velocidad constante $\vec u$.
\end{definition}

\begin{example}
La recta del \cref{ex:g11:vect:twopoints} es también
$\{(1 + 3t,\ 2 + t) : t \in \R\}$: $t = 0$ da $A$ y $t = 1$ da $B$.
Eliminando $t$ ($t = y - 2$, luego $x = 1 + 3(y-2)$) se recupera
$x - 3y + 5 = 0$.
\end{example}

\section{Posiciones relativas de dos rectas}

\begin{proposition}[Paralelas o secantes]\label{prop:g11:vect:positions}
Dos rectas con \omterm{def:g11:vect:direction}{vectores directores} $\vec u$ y $\vec v$ son paralelas si y
solo si $\det(\vec u, \vec v) = 0$. En forma de ecuaciones:
$d : ax + by + c = 0$ y $d' : a'x + b'y + c' = 0$ son paralelas si y solo
si $ab' - a'b = 0$. Dos rectas no paralelas se cortan en exactamente un
punto.
\end{proposition}

\begin{proof}
El paralelismo significa direcciones \omterm{def:g11:vect:collinear}{colineales}, que es el
\cref{thm:g11:vect:collinearity}; con $\vec u\,(-b, a)$ y
$\vec v\,(-b', a')$, el \omterm{def:g11:vect:det}{determinante} es $(-b)a' - (-b')a = ab' - a'b$. Si
las rectas no son paralelas, resolver las dos ecuaciones simultáneamente
(por sustitución o por combinación) lleva a una solución única: el sistema
se comporta como un sistema $2 \times 2$ de \omterm{def:g11:vect:det}{determinante} no nulo.
\end{proof}

\begin{method}[Corte de dos rectas]\label{met:g11:vect:intersection}
Resuelve el sistema de las dos ecuaciones: despeja una incógnita en una
\omterm{def:g10:algebra:equation}{ecuación} y sustitúyela en la otra, o combina las ecuaciones para eliminar
una incógnita. Comprueba siempre el punto obtenido en \emph{las dos}
ecuaciones.
\end{method}

\begin{example}\label{ex:g11:vect:intersection}
Cortemos $d: x - 3y + 5 = 0$ con $d': 2x + y - 4 = 0$. De $d'$:
$y = 4 - 2x$. Sustituyendo en $d$:
\[
x - 3(4 - 2x) + 5 = 0
\iff x - 12 + 6x + 5 = 0
\iff 7x = 7
\iff x = 1,
\]
y entonces $y = 2$. Las rectas se cortan en $(1, 2)$. Comprobación en
$d$: $1 - 6 + 5 = 0$.
\end{example}

\section{Ejercicios}

\begin{exercise}[$\star$]\label{exo:g11:vect:1}
Dados $A(2, -1)$, $B(5, 3)$ y $C(-1, 1)$, calcula las \omterm{def:g10:coordgeom:system}{coordenadas} de
$\vect{AB}$, $\vect{AC}$, $\vect{AB} + \vect{AC}$ y
$2\vect{AB} - 3\vect{AC}$, y el \omterm{prop:g10:coordgeom:midpoint}{punto medio} de $[BC]$.
\end{exercise}

\begin{exercise}[$\star$]\label{exo:g11:vect:2}
¿Son \omterm{def:g11:vect:collinear}{colineales} los \omterm{def:g11:vect:vector}{vectores}?
\[
\vec u\,(4, -6) \text{ y } \vec v\,(-2, 3); \qquad
\vec u\,(2, 5) \text{ y } \vec v\,(3, 7); \qquad
\vec u\,(0, 3) \text{ y } \vec v\,(0, -8).
\]
\end{exercise}

\begin{exercise}[$\star$]\label{exo:g11:vect:3}
Halla una \omterm{def:g10:algebra:equation}{ecuación} general de la recta que pasa por $A(2, 1)$ con \omterm{def:g11:vect:direction}{vector
director} $\vec u\,(3, -1)$, y después la de la recta que pasa por
$B(0, 4)$ y $C(2, 0)$.
\end{exercise}

\begin{exercise}[$\star$]\label{exo:g11:vect:4}
Para la recta $d: 3x - 2y + 6 = 0$, da un \omterm{def:g11:vect:direction}{vector director} y los cortes con
los dos ejes de \omterm{def:g10:coordgeom:system}{coordenadas}, y decide si los puntos $P(2, 6)$ y $Q(1, 4)$
pertenecen a $d$.
\end{exercise}

\begin{exercise}[$\star\star$]\label{exo:g11:vect:5}
¿Están alineados los puntos $A(1, 2)$, $B(4, 8)$, $C(-2, -4)$? La misma
pregunta para $D(0, 1)$, $E(2, 4)$, $F(5, 9)$.
\end{exercise}

\begin{exercise}[$\star\star$]\label{exo:g11:vect:6}
Determina el punto de corte, si existe:
\[
d: 2x + y - 7 = 0 \ \text{ y }\ d': x - y + 1 = 0;
\qquad
e: x - 2y + 3 = 0 \ \text{ y }\ e': -2x + 4y + 1 = 0 .
\]
\end{exercise}

\begin{exercise}[$\star\star$]\label{exo:g11:vect:7}
Da una representación paramétrica de la recta $d: 2x - y + 3 = 0$ y una
\omterm{def:g10:algebra:equation}{ecuación} general de la recta
$\bigl\{(1 + 2t,\ -3 + t) : t \in \R\bigr\}$.
\end{exercise}

\begin{exercise}[$\star\star$]\label{exo:g11:vect:8}
Halla la \omterm{def:g10:algebra:equation}{ecuación} de la recta que pasa por $P(1, -2)$ y es paralela a
$d: 4x - 3y + 2 = 0$, y el valor de $m$ para el que la recta
$mx + 2y - 5 = 0$ es paralela a $d$.
\end{exercise}

\begin{exercise}[$\star\star$]\label{exo:g11:vect:9}
Sean $A(-1, 0)$, $B(3, 2)$ y $C(1, -4)$. Halla una \omterm{def:g10:algebra:equation}{ecuación} general de la
\emph{\omterm{def:g10:stats:median}{mediana}} del triángulo $ABC$ que parte de $C$ (la recta que une $C$
con el \omterm{prop:g10:coordgeom:midpoint}{punto medio} de $[AB]$).
\end{exercise}

\begin{exercise}[$\star\star$]\label{exo:g11:vect:10}
¿Para qué valor o valores del parámetro $m$ son \omterm{def:g11:vect:collinear}{colineales} los \omterm{def:g11:vect:vector}{vectores}
$\vec u\,(m, 2)$ y $\vec v\,(3, m - 1)$?
\end{exercise}

\begin{exercise}[$\star\star\star$]\label{exo:g11:vect:11}
Sea $ABCD$ un paralelogramo, sea $I$ el \omterm{prop:g10:coordgeom:midpoint}{punto medio} de $[AB]$ y sea $J$ el
punto definido por $\vect{DJ} = \frac23 \vect{DI}$. Trabajando en el
\omterm{def:g10:coordgeom:system}{sistema de coordenadas} en el que $A(0,0)$, $B(1,0)$, $D(0,1)$ (de modo que
$C(1,1)$):
\begin{enumerate}
  \item Calcula las \omterm{def:g10:coordgeom:system}{coordenadas} de $I$ y de $J$.
  \item Demuestra que $A$, $J$ y $C$ están alineados. (Un hecho bonito: la
        recta $(DI)$ corta la diagonal $(AC)$ a un tercio de su
        longitud.)
\end{enumerate}
\end{exercise}

\section{Problema: Rumbos de colisión y coordenadas astutas}

\begin{problem}[{Problema de fin de semana --- que dos rumbos se crucen no
es chocar: la máxima aproximación en el mar y teoremas demostrados
eligiendo bien el sistema de referencia}]
\label{pb:g11:vect:1}
Los rumbos rectos de dos barcos se cruzan: ¿tienen que chocar? Solo si
llegan al cruce \emph{al mismo tiempo}: las trayectorias son conjuntos de
puntos, los movimientos son paramétricos, y confundir ambas cosas ha
hundido barcos de verdad. Este problema pone en marcha un pequeño servicio
de tráfico marítimo sobre rectas paramétricas
(\cref{def:g11:vect:parametric}) y pasa después a la geometría pura con el
otro superpoder de las \omterm{def:g10:coordgeom:system}{coordenadas}: \emph{elegirlas} con astucia, como en
el \cref{exo:g11:vect:11}, para que los teoremas clásicos se demuestren
solos.

\medskip
\noindent\textbf{Parte I --- Rectas, de tres maneras.}
\begin{enumerate}
  \item Da una representación paramétrica de la recta que pasa por
        $A(1, 2)$ con \omterm{def:g11:vect:direction}{vector director} $\vec u(3, -1)$ y elimina después el
        parámetro para obtener una \omterm{def:g10:algebra:equation}{ecuación} general.
  \item Para la recta $2x - 5y + 3 = 0$: lee un \omterm{def:g11:vect:direction}{vector director}
        (\cref{thm:g11:vect:cartesian}), halla un punto y escribe una
        representación paramétrica.
  \item Calcula el punto de corte de las dos rectas de las preguntas 1 y 2
        (\cref{met:g11:vect:intersection}).
  \item Confirma con el \omterm{def:g11:vect:det}{determinante} (\cref{thm:g11:vect:collinearity})
        que esas dos rectas tenían que cortarse; y demuestra después que
        $2x - 5y + 3 = 0$ y $-4x + 10y + 1 = 0$ son estrictamente
        paralelas.
  \item ¿Están alineados $A(2, 1)$, $B(5, 3)$, $C(11, 7)$
        (\cref{met:g11:vect:alignment})?
\end{enumerate}

\medskip
\noindent\textbf{Parte II --- El servicio de tráfico marítimo.}
En el instante $t$ (en horas), el barco $P$ está en $(2t,\ 3 + t)$ y el
barco $Q$, en $(10 - t,\ 2t - 1)$ (distancias en millas náuticas).
\begin{enumerate}[resume]
  \item Halla las ecuaciones generales de las dos \emph{trayectorias} y
        calcula dónde se cruzan.
  \item ¿En qué instante pasa cada barco por el punto de cruce? ¿Chocan
        los barcos? Enuncia en una frase la advertencia general:
        trayectorias que se cruzan frente a movimientos que se encuentran.
  \item Calcula la distancia al cuadrado $D^2(t)$ entre los barcos en el
        instante $t$ y redúcela a un polinomio de segundo grado en $t$.
        ¿En qué instante están más cerca y a qué distancia llegan a estar
        (la tecnología del vértice del \cref{pb:g11:quad:1})?
  \item El reglamento exige una separación mínima de $1$ milla. ¿Se
        incumple? Si es así, resuelve $D^2(t) < 1$ y da la duración del
        incumplimiento.
  \item Explica por qué, para \emph{dos} barcos cualesquiera con rumbo
        recto y velocidad constante, $D^2(t)$ es siempre una \omterm{def:g11:quad:quadratic}{función de
        segundo grado} de $t$, de manera que el ``punto de máxima
        \omterm{def:g10:numbers:approx}{aproximación}'' es siempre un cálculo de vértice.
  \item Intercepción: una patrullera parte del origen en $t = 0$ con
        velocidad constante $(a, b)$ y debe encontrarse con el barco $P$
        exactamente en $t = 3$. Calcula el $(a, b)$ necesario y la
        velocidad de la patrullera.
\end{enumerate}

\medskip
\noindent\textbf{Parte III --- \omterm{def:g10:coordgeom:system}{Coordenadas} astutas.}
Trabajamos en el sistema del \cref{exo:g11:vect:11}: el paralelogramo
$ABCD$ pasa a ser $A(0,0)$, $B(1,0)$, $C(1,1)$, $D(0,1)$; sea $I$ el \omterm{prop:g10:coordgeom:midpoint}{punto
medio} de $[AB]$ y $K$ el \omterm{prop:g10:coordgeom:midpoint}{punto medio} de $[CD]$.
\begin{enumerate}[resume]
  \item Calcula una \omterm{def:g10:algebra:equation}{ecuación} general de la recta $(DI)$.
  \item Corta $(DI)$ con la diagonal $(AC)$: recupera el hecho del
        \cref{exo:g11:vect:11}, o sea, que el corte está exactamente a un
        tercio del \omterm{def:g10:stats:quartiles}{recorrido} de la diagonal.
  \item Ahora la recta $(BK)$: halla su \omterm{def:g10:algebra:equation}{ecuación}, córtala con $(AC)$ y
        concluye el teorema clásico completo: \emph{las cevianas $(DI)$ y
        $(BK)$ dividen la diagonal $(AC)$ en tres partes iguales}.
  \item ¿Por qué era legítimo demostrar el teorema en ese único sistema
        particular? (¿Qué conserva la elección de \omterm{def:g10:coordgeom:system}{coordenadas} y de qué
        habla el enunciado?) Responde en dos o tres frases.
  \item Otra trisección para practicar: sea $E$ el \omterm{prop:g10:coordgeom:midpoint}{punto medio} de $[BC]$.
        ¿Dónde corta la recta $(AE)$ a la \emph{otra} diagonal $(DB)$?
\end{enumerate}

\medskip
\noindent\textbf{Parte IV --- Acertijos de posición.}
\begin{enumerate}[resume]
  \item ¿Son concurrentes las tres rectas $x + y = 4$, \ $2x - y = 2$, \
        $x - 2y = -2$? (Corta dos y prueba con la tercera.)
  \item Para cada real $m$, sea $d_m$ la recta $mx - y + 2 - 3m = 0$.
        Demuestra que \emph{todas} las $d_m$ pasan por un mismo \omterm{pb:g10:functions:1}{punto
        fijo}. (Agrupa los términos en $m$.)
  \item En la familia $(d_m)$: ¿qué miembro es paralelo a $y = 2x + 7$?
        ¿Cuál pasa por el origen?
  \item Final: los tres disfraces de una recta, el general (pertenencia
        resuelta con una sustitución), el explícito (\omterm{def:g10:reffunc:affine}{pendiente} y \omterm{def:g10:functions:graph}{gráfica} a
        la vista) y el paramétrico (movimiento, horarios, intercepción);
        el \omterm{def:g11:vect:det}{determinante} como oráculo del paralelismo; y el sistema de
        referencia bien elegido como fábrica de teoremas. Una frase para
        cada uno.
\end{enumerate}
\end{problem}
